% Part: applied-modal-logics
% Chapter: epistemic-logic
% Section: relational-models

\documentclass[../../../include/open-logic-section]{subfiles}

\begin{document}

\olfileid{aml}{el}{rel}

\olsection{সম্পর্কভিত্তিক মডেল}

জ্ঞানতাত্ত্বিক যুক্তিবিদ্যার মৌলিক অর্থতাত্ত্বিক ধারণাটি সাধারণ মোডাল যুক্তিবিদ্যার মতোই। সম্পর্কভিত্তিক মডেলে এখনও জগৎগুলির একটি সমষ্টি এবং কোন জগতে কোন !!{propositional
variable}s ``সত্য'' তা নির্ধারণকারী একটি নির্ধারণ থাকে। আর কেবল একজন কর্তার কথা হলে, যথারীতি একটি অভিগম্যতা-সম্পর্ক থাকে। কিন্তু বহু-কর্তা জ্ঞানতাত্ত্বিক যুক্তিবিদ্যায় সেই একটি অভিগম্যতা-সম্পর্কের বদলে কয়েকটি সম্পর্কের সমষ্টি থাকে—কর্তা-চিহ্নের সমষ্টি~$G$-এর প্রত্যেক~$a$-র জন্য একটি করে।

একটি \emph{সম্পর্কভিত্তিক মডেল}-এ থাকে জগৎগুলির একটি সমষ্টি, প্রত্যেক কর্তার জন্য একটি করে জগৎগুলির মধ্যকার দ্বিপদী অভিগম্যতা-সম্পর্ক, এবং কোন জগতে কোন !!{propositional variable}s সত্য তা নির্ধারণকারী একটি নির্ধারণ।

\begin{defn}
  বহু-কর্তা জ্ঞানতাত্ত্বিক ভাষার একটি \emph{মডেল} হল ত্রয়ী
  $\mModel{M} = \tuple{W, R, V}$, যেখানে
  \begin{enumerate}
  \item $W$ ``জগৎ''-গুলির একটি অশূন্য সমষ্টি,
  \item প্রত্যেক $a \in G$-এর জন্য ${R}_a$ হল~$W$-এর ওপর একটি দ্বিপদী অভিগম্যতা-সম্পর্ক, এবং
  \item $V$ এমন একটি অপেক্ষক, যা প্রত্যেক !!{propositional
    variable}~$p$-কে সম্ভাব্য জগৎগুলির একটি সমষ্টি $V(p)$ নির্ধারণ করে দেয়।
  \end{enumerate}
  $R_a ww'$ হলে আমরা বলি, $w$~থেকে $w'$ \emph{a-র কাছে অভিগম্য}। $w \in V(p)$ হলে বলি, $p$ জগৎ~$w$-তে \emph{সত্য}।
\end{defn}

পদ্ধতিটি স্বাভাবিক মোডাল যুক্তিবিদ্যার মতোই, শুধু অতিরিক্ত অভিগম্যতা-সম্পর্ক যোগ হয়েছে। নির্দিষ্ট কর্তার অভিগম্যতা-সম্পর্ককে সাধারণত তার তথ্যাবস্থা সম্পর্কে কিছু বোঝাতে ব্যবহার করব। যেমন, $R_a ww'$-কে প্রায়ই এই অর্থে নিই যে জগৎ~$w$-তে $a$-র কাছে থাকা তথ্যের সঙ্গে~$w'$ সঙ্গতিপূর্ণ। অন্যভাবে বললে, $w$-তে সে জগৎ~$w$ ও জগৎ~$w'$-এর পার্থক্য বুঝতে পারে না।

\end{document}
